\section{INTRODUCTION}

Python has witnessed remarkable growth over the past decade and has become one
of the most widely used programming languages worldwide.
Since 2018, its popularity has surged, and it has been ranked as the most
popular programming language according to multiple surveys since 2022.
This rise can be attributed to Python’s versatility and applicability across a
broad spectrum of domains, including web development, networking, data
analytics, scientific computing, automation, and, most prominently, artificial
intelligence and machine learning.
In particular, the rapid proliferation of AI-related applications, such as deep
learning frameworks (e.g., TensorFlow, PyTorch) and natural language processing
libraries (e.g., Hugging Face Transformers), has further accelerated Python’s
adoption among both industry professionals and researchers.

In addition to its broad applicability, Python’s concise and
natural-language-like syntax lowers the barrier to entry for new developers.
Features such as dynamic typing, high-level data structures, and rich standard
libraries allow programmers to express complex ideas with minimal boilerplate
code.
As a result, Python has become a first-choice language for prototyping,
experimentation, and educational purposes, fostering a vast ecosystem of
open-source libraries and active developer communities.

However, Python’s dynamic and interpreted nature introduces significant
challenges for program reliability.
Unlike compiled languages, where many errors can be caught at compile time,
Python programs often surface errors only during execution.
Among these runtime errors, type-related issues such as \texttt{TypeError} and
\texttt{AttributeError} are particularly prevalent, accounting for a
substantial portion of all reported bugs—studies suggest approximately 30\% of
runtime errors in Python stem from type mismatches.
This prevalence of type errors not only affects software reliability but also
complicates maintenance and refactoring of large codebases, where inconsistent
or unexpected type usage can propagate subtle bugs across modules.

Consequently, type checking and type analysis have become critical for
improving the reliability, maintainability, and robustness of Python programs.
The introduction of optional type annotations in Python 3.5 provided a
mechanism for developers to express expected types explicitly, enabling static
verification while preserving the language’s inherent flexibility.
Tools such as \texttt{mypy}, \texttt{Pyre}, and \texttt{Pyright} leverage these
annotations to detect type inconsistencies before execution, reducing the
likelihood of runtime failures.
Beyond static checking, hybrid approaches combining static and runtime
monitoring, as well as emerging AI-based techniques, have begun to address
Python’s dynamic and flexible behavior more effectively.

Despite these advancements, Python type analysis remains challenging due to
several unique characteristics of the language.
First, Python follows a form of structural typing, often referred to as “duck
typing,” where object compatibility is determined by available attributes and
methods rather than nominal type hierarchies.
Second, variable types in Python can change dynamically depending on the
control flow, requiring type inference mechanisms to be flow-sensitive to
achieve accuracy.
Third, Python supports advanced dynamic features, including reflection, runtime
code generation through \texttt{eval} and \texttt{exec}, and dynamic module
loading, all of which complicate static reasoning.
Finally, Python data structures such as lists, dictionaries, and sets
frequently contain heterogeneous elements whose types may vary over time,
making traditional uniform-type assumptions insufficient.
These language characteristics collectively create a significant gap between
static type assumptions and the actual behavior of programs, motivating the
development of specialized analysis techniques.

In this context, the present study aims to provide a comprehensive survey of
Python type analysis techniques and tools, categorizing them according to their
underlying methodologies and examining their strengths and limitations.
Specifically, this paper addresses the following research questions (RQs):

\begin{itemize}
  \item \textbf{RQ1} : How can Python’s dynamic nature be
    preserved while ensuring program stability through type checking?
  \item \textbf{RQ2} : What are the main challenges in type inference and type
    analysis for Python programs?
  \item \textbf{RQ3} : What technical characteristics and limitations define
    existing Python type analysis tools?
\end{itemize}

To answer these questions, we analyze Python type analysis tools spanning
several paradigms, including static analyzers leveraging type annotations,
runtime and hybrid monitoring systems, constraint-solving and symbolic
execution-based approaches, and recently developed AI-powered methods enabled
by large language models.
By systematically examining these tools and approaches, we aim to clarify the
current landscape of Python type analysis, highlight open challenges, and
identify promising directions for future research.

In summary, this paper contributes a structured overview of Python type
analysis, emphasizing the interplay between dynamic flexibility and the need
for program reliability, and provides insights for both tool developers and
researchers aiming to advance Python’s type-checking ecosystem.

